\subsection{Restriction of the base field}

In this number, one is given two complete non-discrete valued commutative fields K and L, as well as an isomorphism \(\sigma\) of the valued field K onto a subfield of L. If E is a Banach space over L, one denotes by \(\sigma_{*}(\mathbf{E})\) the vector space over K obtained by restriction of scalars (cf.~Alg., chap.~II, 3rd ed., § 8); the given topology on E is compatible with the structure of K-vector space, and \(\sigma_{*}(\mathbf{E})\) is a Banach space over K.

5.14.1. Let X be an analytic manifold over L and let \(c = (U, \varphi, E)\) be a chart of X. The map \(\varphi\) is a bijection of U onto an open subset of \(\sigma_{*}(E)\) and the triple \(c_{\sigma} = (U, \varphi, \sigma_{*}(E))\) is a chart of X. There exists on X a structure of analytic manifold over K and only one for which \(c_{\sigma}\) is a chart for every chart c of the L-analytic manifold X. One denotes by \(X_{\sigma}\) the analytic manifold over K thus obtained and one says that \(X_{\sigma}\) is obtained from X by restriction of scalars (from L to K by means of \(\sigma\) ). The topological space underlying \(X_{\sigma}\) is the same as that underlying X.

5.14.2. Examples:

\begin{enumerate}
\def\labelenumi{(\alph{enumi})}
\item
  One takes K = R, L = C, σ being the canonical injection of R into C. Thus, every complex analytic manifold is canonically equipped with a structure of real analytic manifold; this real analytic structure itself defines differential structures of class \(C^{r}\) for every r.
\item
  One takes K = C, L = C, σ being the conjugation \(x \mapsto \bar{x}\) . One thus associates to every complex analytic manifold X a complex analytic manifold \(\bar{X}\) , called the conjugate of X. If f is a complex-valued function, defined on an open subset U of X, f is analytic for the structure of \(\bar{X}\) if and only if the conjugate function \(\bar{f}\) is analytic for the structure of X. The manifolds X and \(\bar{X}\) define by restriction of scalars the same real analytic manifold.
\end{enumerate}

5.14.3. Let X be an analytic manifold over K, Y an analytic manifold over L. A map \(f:X\to Y\) is said to be \(\sigma\) -analytic if it is a morphism of X into \(Y_{\sigma}\) . If \(K\subset L\) , \(\sigma\) is the injection of K into L, and X is an analytic manifold over L, one calls K-analytic map a map \(\sigma\) -analytic from \(X_{\sigma}\) into Y.

5.14.4. Let V be a Banach space over L. One has \(V_{\sigma} = \sigma_{*}(V)\) ; the canonical analytic variety structure over K of \(\sigma_{*}(V)\) is deduced by restriction of scalars from the canonical analytic variety structure over L of V.

5.14.5. Let X be an analytic variety over L, and let \(a \in X\) . Let c be a chart of X at a; then \(c_{\sigma}\) is a chart of \(X_{\sigma}\) at a, and one deduces from it isomorphisms

\[
\theta_ {c} \colon \mathrm{E} \rightarrow \mathrm{T} _ {a} (\mathrm{X}), \quad \theta_ {c _ {\sigma}} \colon \sigma_ {*} (\mathrm{E}) \rightarrow \mathrm{T} _ {a} (\mathrm{X} _ {\sigma})
\]

whence an isomorphism \(\theta_{c_{\sigma}} \circ \sigma_{*}(\theta_{c})^{-1}\) of \(\sigma_{*}(\mathrm{T}_{a}(\mathrm{X}))\) onto \(\mathrm{T}_{a}(\mathrm{X}_{\sigma})\) ; this isomorphism does not depend on the choice of \(c\) ; it allows one to identify \(\mathrm{T}_{a}(\mathrm{X}_{\sigma})\) with \(\sigma_{*}(\mathrm{T}_{a}(\mathrm{X}))\) and even with \(\mathrm{T}_{a}(\mathrm{X})\) by abuse of notation.

If \(f\) is an analytic map on L from X into a manifold Y, the tangent linear map to \(f\) at \(a\) ( \(f\) being regarded as a morphism of \(\mathrm{X}_{\sigma}\) to \(\mathrm{Y}_{\sigma}\) ) is equal to \(\sigma_{*}(\mathrm{T}_{a}(f))\) .

5.14.6. Let X and Y be two analytic varieties over L, and let \(f: \mathbf{X}_{\sigma} \to \mathbf{Y}\) be a continuous \(\sigma\) -analytic map. Suppose that the characteristic of K is 0. Then, for \(f\) to be analytic over L, it is necessary and sufficient that, for every \(a \in \mathbf{X}\) , the map \(\mathrm{T}_{a}(f)\) be L-linear.

When K = R, L = C (case (a) of no. 5.14.2), one has a more precise result: if X and Y are finite-dimensional and if \(f: X \to Y\) is a map of class \(C^{1}\) whose tangent map at every point of X is C-linear, then \(f\) is complex analytic.

5.14.7. Let X be a complex analytic manifold and g a map from X into itself satisfying the conditions:

\begin{enumerate}
\def\labelenumi{(\roman{enumi})}
\item
  One has \(g \circ g = \mathrm{Id}_{\mathbf{X}}\) .
\item
  The map \(g\) is an isomorphism of the analytic variety \(X\) on the conjugate variety \(\overline{X}\) (5.14.2).
\end{enumerate}

The set \(X_{0}\) of points x of X such that \(g(x) = x\) is a closed analytic submanifold of the real analytic manifold underlying X. For \(x \in X_{0}\) , one has \(\mathbf{T}_{x}(\mathbf{X}) = \mathbf{T}_{x}(\mathbf{X}_{0}) \oplus i\mathbf{T}_{x}(\mathbf{X}_{0})\) .

Let U be a connected open subset of X, and let f and g be two complex analytic maps from U into a complex separated locally convex space or into a complex separated analytic manifold. If f and g coincide on a nonempty subset of \(U \cap X_{0}\) , open in \(X_{0}\) , then f = g.

Suppose \(X_{0}\) is paracompact. If f is a real analytic map from \(X_{0}\) into a separated locally convex space or into a separated complex analytic manifold, there exists an open neighborhood U of \(X_{0}\) in X and a complex analytic map from U into the space of values of f, extending f. Any two such extensions coincide on a neighborhood of \(X_{0}\) in X.

Suppose X is of finite dimension. For every point \(a \in X_{0}\) , there exists a system of (complex) coordinates \(\zeta^{1}, \ldots, \zeta^{n}\) in an open neighborhood U of \(a\) , such that \(\zeta^{i} \circ g = \bar{\zeta}^{i}\) for \(1 \leqslant i \leqslant n\) ; the restriction \(\xi^{i}\) of \(\zeta^{i}\) to \(U \cap X_{0}\) then takes real values and \((\xi^{1}, \ldots, \xi^{n})\) is a system of coordinates at \(a\) of the real analytic manifold \(X_{0}\) .

5.14.8. For every paracompact real analytic manifold Y, there exists a pair \((X, g)\) consisting of a complex analytic manifold X and a map g from X into itself satisfying conditions (i) and (ii) of 5.14.7, and an isomorphism f from Y onto \(X_{0}\) . We say that X (endowed with f and g) is a complexification of Y.

